\sectionkicker{Source-backed technical notes}
\subsection*{能力密度・Serving・Localityを再設計する: Technical Notes}
本欄は記事本文で圧縮した一次資料上の情報を、比較・再検証しやすい形で整理したものである。時系列、確認済みの主張、留意点、一次資料URLを掲載する。
\smallskip
\noindent{\small\textit{Technical Notesでは、一次資料から確認した公開・提供・研究上の事実を記載する。数値・能力評価や提供元・著者による比較表現は、独立再現済みの結果ではなく帰属claimとして扱う。}}\par
\medskip
\subsection*{Theme at a glance}
\addcontentsline{toc}{subsection}{Theme at a glance}
\begin{center}
\footnotesize
\begin{tabularx}{\linewidth}{@{}p{0.20\linewidth}p{0.10\linewidth}p{0.14\linewidth}X@{}}
\toprule
資料 & 位置づけ & 種別 & 時系列 \\
\midrule
FlashAttention-3 & 主要資料 & Framework & 2024-07-11             \\
Prompt Caching in the API & 主要資料 & API & 2024-10-01             \\
Model Distillation in the API & 主要資料 & API & 2024-10-01             \\
LTX-Video & 主要資料 & オープンウェイト & 2024-11-21             \\
MagicDec: Breaking the Latency-Throughput Tradeoff for Long Context Generation with Speculative Decoding & 補足資料 & 論文 & 2024-08-20             \\
Qwen2.5-Turbo 1M context & 補足資料 & モデル更新 & 2024-11-15             \\
Quantized Llama 3.2 lightweight models & 補足資料 & モデル更新 & 2024-10-24             \\
\bottomrule
\end{tabularx}
\end{center}
\normalsize


\medskip

\begin{technicalnote}{FlashAttention-3}{主要資料}
\begingroup
% reader-facing Technical Notes break policy
\widowpenalty=10000
\clubpenalty=10000
\displaywidowpenalty=10000
\begin{tabularx}{\linewidth}{@{}>{\bfseries}p{0.22\linewidth}X@{}}
組織 & FlashAttention authors / PyTorch \\
種別 & Framework \\
時系列 & 2024-07-11 \\
\end{tabularx}
\smallskip
{\bfseries 一次資料から整理したtechnical points}
\begin{itemize}[leftmargin=1.5em,itemsep=0.35em]
\item \textbf{一次情報で確認できる事実}: FlashAttention authors / PyTorchの一次資料により、FlashAttention-3について2024-07-11にSoftware公開を確認した。 対象event近傍の一次資料から FlashAttention-3 Hopper/TMA/WGMMA optimization / Hopper TMA / WGMMA活用 / FP8 を確認できる。
\end{itemize}
\begin{minipage}{\linewidth}
% reader-facing Technical Notes generic boundary/source tail group
\begin{samepage}
{\bfseries 一次資料}
\begingroup\sloppy
\Urlmuskip=0mu plus 2mu\relax
\begin{itemize}[leftmargin=1.5em,itemsep=0.25em]
\item {\scriptsize\url{https://pytorch.org/blog/flashattention-3/}}
\end{itemize}
\endgroup
\end{samepage}
\end{minipage}
\endgroup
\end{technicalnote}
\medskip

\begin{technicalnote}{Prompt Caching in the API}{主要資料}
\begingroup
% reader-facing Technical Notes break policy
\widowpenalty=10000
\clubpenalty=10000
\displaywidowpenalty=10000
\begin{tabularx}{\linewidth}{@{}>{\bfseries}p{0.22\linewidth}X@{}}
組織 & OpenAI \\
種別 & API \\
時系列 & 2024-10-01 \\
\end{tabularx}
\smallskip
{\bfseries 一次資料から整理したtechnical points}
\begin{itemize}[leftmargin=1.5em,itemsep=0.35em]
\item \textbf{一次情報で確認できる事実}: OpenAIの一次資料により、Prompt Caching in the APIについて2024-10-01にAPI提供開始を確認した。 対象event近傍の一次資料から recent-input reuse / caching discount を確認できる。
\end{itemize}
\begin{minipage}{\linewidth}
% reader-facing Technical Notes generic boundary/source tail group
\begin{samepage}
{\bfseries 一次資料}
\begingroup\sloppy
\Urlmuskip=0mu plus 2mu\relax
\begin{itemize}[leftmargin=1.5em,itemsep=0.25em]
\item {\scriptsize\url{https://openai.com/index/api-prompt-caching}}
\end{itemize}
\endgroup
\end{samepage}
\end{minipage}
\endgroup
\end{technicalnote}
\medskip

\begin{technicalnote}{Model Distillation in the API}{主要資料}
\begingroup
% reader-facing Technical Notes break policy
\widowpenalty=10000
\clubpenalty=10000
\displaywidowpenalty=10000
\begin{tabularx}{\linewidth}{@{}>{\bfseries}p{0.22\linewidth}X@{}}
組織 & OpenAI \\
種別 & API \\
時系列 & 2024-10-01 \\
\end{tabularx}
\smallskip
{\bfseries 一次資料から整理したtechnical points}
\begin{itemize}[leftmargin=1.5em,itemsep=0.35em]
\item \textbf{一次情報で確認できる事実}: OpenAIの一次資料により、Model Distillation in the APIについて2024-10-01にAPI提供開始を確認した。 対象event近傍の一次資料から fine-tuning from frontier-model outputs / small model / cost-efficient deployment を確認できる。
\end{itemize}
\begin{minipage}{\linewidth}
% reader-facing Technical Notes generic boundary/source tail group
\begin{samepage}
{\bfseries 一次資料}
\begingroup\sloppy
\Urlmuskip=0mu plus 2mu\relax
\begin{itemize}[leftmargin=1.5em,itemsep=0.25em]
\item {\scriptsize\url{https://openai.com/index/api-model-distillation}}
\end{itemize}
\endgroup
\end{samepage}
\end{minipage}
\endgroup
\end{technicalnote}
\medskip

\begin{technicalnote}{LTX-Video}{主要資料}
\begingroup
% reader-facing Technical Notes break policy
\widowpenalty=10000
\clubpenalty=10000
\displaywidowpenalty=10000
\begin{tabularx}{\linewidth}{@{}>{\bfseries}p{0.22\linewidth}X@{}}
組織 & Lightricks \\
種別 & オープンウェイト \\
時系列 & 2024-11-21 \\
\end{tabularx}
\smallskip
{\bfseries 一次資料から整理したtechnical points}
\begin{itemize}[leftmargin=1.5em,itemsep=0.35em]
\item \textbf{一次情報で確認できる事実}: Lightricksの一次資料により、LTX-Videoについて2024-11-21にモデル公開を確認した。 対象event近傍の一次資料から text-to-video / image-to-video を確認できる。
\end{itemize}
\begin{minipage}{\linewidth}
% reader-facing Technical Notes generic boundary/source tail group
\begin{samepage}
{\bfseries 一次資料}
\begingroup\sloppy
\Urlmuskip=0mu plus 2mu\relax
\begin{itemize}[leftmargin=1.5em,itemsep=0.25em]
\item {\scriptsize\url{https://github.com/Lightricks/LTX-Video/blob/main/README.md}}
\end{itemize}
\endgroup
\end{samepage}
\end{minipage}
\endgroup
\end{technicalnote}
\medskip

\begin{technicalnote}{MagicDec: Breaking the Latency-Throughput Tradeoff for Long Context Generation with Speculative Decoding}{補足資料}
\begingroup
% reader-facing Technical Notes break policy
\widowpenalty=10000
\clubpenalty=10000
\displaywidowpenalty=10000
\begin{tabularx}{\linewidth}{@{}>{\bfseries}p{0.22\linewidth}X@{}}
組織 & - \\
種別 & 論文 \\
時系列 & 2024-08-20 \\
\end{tabularx}
\smallskip
{\bfseries 一次資料から整理したtechnical points}
\begin{itemize}[leftmargin=1.5em,itemsep=0.35em]
\item \textbf{一次情報で確認できる事実}: 提供元の一次資料により、MagicDec: Breaking the Latency-Throughput Tradeoff for Long Context Generation with Speculative Decodingについて2024-08-20にPaper First Submissionを確認した。 対象event近傍の一次資料から 8B parameter scale / speculative decoding / KV-cache を確認できる。
\end{itemize}
\begin{minipage}{\linewidth}
% reader-facing Technical Notes generic boundary/source tail group
\begin{samepage}
{\bfseries 一次資料}
\begingroup\sloppy
\Urlmuskip=0mu plus 2mu\relax
\begin{itemize}[leftmargin=1.5em,itemsep=0.25em]
\item {\scriptsize\url{https://arxiv.org/abs/2408.11049}}
\end{itemize}
\endgroup
\end{samepage}
\end{minipage}
\endgroup
\end{technicalnote}
\medskip

\begin{technicalnote}{Qwen2.5-Turbo 1M context}{補足資料}
\begingroup
% reader-facing Technical Notes break policy
\widowpenalty=10000
\clubpenalty=10000
\displaywidowpenalty=10000
\begin{tabularx}{\linewidth}{@{}>{\bfseries}p{0.22\linewidth}X@{}}
組織 & Qwen \\
種別 & モデル更新 \\
時系列 & 2024-11-15 \\
\end{tabularx}
\smallskip
{\bfseries 一次資料から整理したtechnical points}
\begin{itemize}[leftmargin=1.5em,itemsep=0.35em]
\item \textbf{一次情報で確認できる事実}: Qwenの一次資料により、Qwen2.5-Turbo 1M contextについて2024-11-15にモデル公開を確認した。 対象event近傍の一次資料から 9B parameter scale / 128K context / 1M context を確認できる。
\end{itemize}
\begin{minipage}{\linewidth}
% reader-facing Technical Notes generic boundary/source tail group
\begin{samepage}
{\bfseries 一次資料}
\begingroup\sloppy
\Urlmuskip=0mu plus 2mu\relax
\begin{itemize}[leftmargin=1.5em,itemsep=0.25em]
\item {\scriptsize\url{https://qwenlm.github.io/blog/qwen2.5-turbo/}}
\end{itemize}
\endgroup
\end{samepage}
\end{minipage}
\endgroup
\end{technicalnote}
\medskip

\begin{technicalnote}{Quantized Llama 3.2 lightweight models}{補足資料}
\begingroup
% reader-facing Technical Notes break policy
\widowpenalty=10000
\clubpenalty=10000
\displaywidowpenalty=10000
\begin{tabularx}{\linewidth}{@{}>{\bfseries}p{0.22\linewidth}X@{}}
組織 & Meta \\
種別 & モデル更新 \\
時系列 & 2024-10-24 \\
\end{tabularx}
\smallskip
{\bfseries 一次資料から整理したtechnical points}
\begin{itemize}[leftmargin=1.5em,itemsep=0.35em]
\item \textbf{一次情報で確認できる事実}: Metaの一次資料により、Quantized Llama 3.2 lightweight modelsについて2024-10-24にモデル公開を確認した。 対象event近傍の一次資料から 1B parameter scale / 3B parameter scale / 8K context / 8-bit quantization / Quantization-Aware Training (QAT) / SpinQuant / LoRA を確認できる。
\end{itemize}
\begin{minipage}{\linewidth}
% reader-facing Technical Notes generic boundary/source tail group
\begin{samepage}
{\bfseries 一次資料}
\begingroup\sloppy
\Urlmuskip=0mu plus 2mu\relax
\begin{itemize}[leftmargin=1.5em,itemsep=0.25em]
\item {\scriptsize\url{https://ai.meta.com/blog/meta-llama-quantized-lightweight-models/}}
\end{itemize}
\endgroup
\end{samepage}
\end{minipage}
\endgroup
\end{technicalnote}
\medskip
